% ----- CHAUPAI 27 -----

\newpage

\subsection*{\deva{सप्तविंशतितम चौपाई} (Twenty-Seventh Chaupai)}
\addcontentsline{toc}{subsection}{\deva{सप्तविंशतितम चौपाई} (Twenty-Seventh Chaupai) — \deva{सब पर...}}

\begin{minipage}[t]{0.55\textwidth}
\vspace{0pt}

{\large \linespread{1.5}\selectfont
\deva{सब पर राम तपस्वी राजा।}\\
\deva{तिन के काज सकल तुम साजा॥}
\par}

\vspace{0.5em}

{\large \linespread{1.5}\selectfont
\telu{సబ పర రామ తపస్వీ రాజా।}\\
\telu{తిన కే కాజ సకల తుమ సాజా॥}
\par}

\vspace{0.5em}

{\large \linespread{1.3}\selectfont
\textit{saba para rāma tapasvī rājā |}\\
\textit{tina ke kāja sakala tuma sājā ||}
\par}
\end{minipage}%
\hfill
\begin{minipage}[t]{0.42\textwidth}
\vspace{0pt}
\begin{tcolorbox}[anchorbox]
\textbf{Executing the Vision:} Even the greatest kings and visionaries rely on dedicated individuals to execute the hardest tasks on the ground. A grand vision or a brilliant idea is meaningless without a hard-working team willing to put in the effort to make it a reality.
\vspace{0.5em}\newline Even the greatest visionary leaders required his relentless, on-the-ground execution to turn their goals into reality.\newline\null\hfill\href{https://www.valmiki.iitk.ac.in/sloka?field_kanda_tid=5&language=dv&field_sarga_value=65}{\textit{Sundara Kanda, 65:26}}
\end{tcolorbox}
\end{minipage}

\vspace{0.5em}
\subsubsection*{\textbf{Visual Rhythm Map:}}
\begin{table}[h!]
\centering
\renewcommand{\arraystretch}{1.2}
\setlength{\tabcolsep}{0pt}
\begin{tabularx}{\textwidth}{|*{16}{>{\centering\arraybackslash\hsize=1.0\hsize}X|}}
\hline
\multicolumn{2}{|>{\centering\arraybackslash\hsize=2.0\hsize}X|}{\textbf{\laghu{स}\laghu{ब}} \par (2)} &
\multicolumn{2}{>{\centering\arraybackslash\hsize=2.0\hsize}X|}{\textbf{\laghu{प}\laghu{र}} \par (2)} &
\multicolumn{3}{>{\centering\arraybackslash\hsize=3.0\hsize}X|}{\textbf{\guru{रा}\laghu{म}} \par (3)} &
\multicolumn{5}{>{\centering\arraybackslash\hsize=5.0\hsize}X|}{\textbf{\laghu{त}\guru{प}\guru{स्वी}} \par (5)} &
\multicolumn{4}{>{\centering\arraybackslash\hsize=4.0\hsize}X|}{\textbf{\guru{रा}\guru{जा}} \par (4)} \\
\hline
\end{tabularx}
\vspace{-1.5em}
\end{table}

\begin{table}[h!]
\centering
\renewcommand{\arraystretch}{1.2}
\setlength{\tabcolsep}{0pt}
\begin{tabularx}{\textwidth}{|*{16}{>{\centering\arraybackslash\hsize=1.0\hsize}X|}}
\hline
\multicolumn{2}{|>{\centering\arraybackslash\hsize=2.0\hsize}X|}{\textbf{\laghu{ति}\laghu{न}} \par (2)} &
\multicolumn{2}{>{\centering\arraybackslash\hsize=2.0\hsize}X|}{\textbf{\guru{के}} \par (2)} &
\multicolumn{3}{>{\centering\arraybackslash\hsize=3.0\hsize}X|}{\textbf{\guru{का}\laghu{ज}} \par (3)} &
\multicolumn{3}{>{\centering\arraybackslash\hsize=3.0\hsize}X|}{\textbf{\laghu{स}\laghu{क}\laghu{ल}} \par (3)} &
\multicolumn{2}{>{\centering\arraybackslash\hsize=2.0\hsize}X|}{\textbf{\laghu{तु}\laghu{म}} \par (2)} &
\multicolumn{4}{>{\centering\arraybackslash\hsize=4.0\hsize}X|}{\textbf{\guru{सा}\guru{जा}} \par (4)} \\
\hline
\end{tabularx}
\vspace{-1.5em}
\end{table}

\vspace{1em}
\textbf{ \deva{सरल-संस्कृतम्}:}\\
 \deva{सर्वेषामुपरि तपस्वी राजा श्रीरामः अस्ति।}\\
 \deva{तस्य अपि सर्वाणि कार्याणि भवान् एव सिद्धानि कृतवान् (साधितवान्)॥}\\


Above all beings is Lord Rama, the ascetic King.\\
Yet, you are the one who accomplished all of even His tasks!


\vspace{1em}
\newpage
\textbf{\deva{प्रतिपदार्थः} (Meaning of each word):}
\begin{table}[h!]
\centering
\renewcommand{\arraystretch}{1.2}
\begin{tabularx}{\textwidth}{@{} l l X X @{}}
\toprule
\textbf{Awadhi} & \textbf{Sanskrit} & \textbf{English} & \textbf{Grammar \& Notes} \\
\midrule
\deva{सब पर} / \telu{సబ పర} / \textit{saba para}  &  \deva{सर्वेषाम् उपरि}  &  Above everyone  &   \\
\deva{राम} / \telu{రామ} / \textit{rāma}  &  \deva{रामः}  &  Rama  &   \\
\deva{तपस्वी राजा} / \telu{తపస్వీ రాజా} / \textit{tapasvī rājā}  &  \deva{तपस्वी राजा}  &  Ascetic King  & Rama ruled as an ascetic \\
\deva{तिन के} / \telu{తిన కే} / \textit{tina ke}  &  \deva{तस्य अपि}  &  Even His  & Honorific plural genitive \\
\deva{काज} / \telu{కాజ} / \textit{kāja}  &  \deva{कार्याणि}  &  Tasks / Work  & (from Kaarya) \\ 
\deva{सकल} / \telu{సకల} / \textit{sakala}  &  \deva{सकलानि / सर्वाणि}  &  All / Entire  &   \\
\deva{तुम} / \telu{తుమ} / \textit{tuma}  &  \deva{त्वम्}  &  You  &   \\
\deva{साजा} / \telu{సాజా} / \textit{sājā}  &  \deva{साधितवान्}  &  Accomplished / Arranged  &   \\
\bottomrule
\end{tabularx}
\end{table}



\newpage
